\documentclass[12pt]{article}

\usepackage[T1]{fontenc}
\usepackage[utf8]{inputenc}
\usepackage{amsmath,amssymb,amsthm,enumitem}
\usepackage[margin=1in]{geometry}
\usepackage[colorlinks=true,linkcolor=blue]{hyperref}

\setlength{\parindent}{0pt}
\setlength{\parskip}{5pt}
\setlength{\abovedisplayskip}{7pt}
\setlength{\belowdisplayskip}{7pt}
\setlist{leftmargin=1.5em,itemsep=4pt,topsep=4pt}

\newcommand{\op}{\diamond}
\newcommand{\fin}{\mathrm{fin}}
\newcommand{\E}[1]{\mathrm{E#1}}

\theoremstyle{definition}
\newtheorem{problem}{Problem}
\newtheorem*{hint}{Hint}
\newtheorem*{goal}{Final goal}

\theoremstyle{remark}
\newtheorem*{warning}{Warning}

\title{A Finite Magma Challenge for Equation 677}
\author{}
\date{May 2026}

\begin{document}
\maketitle

\begin{problem}[Main question]
Let $M$ be a finite nonempty set with a binary operation $\op$.  Suppose that, for all
$x,y\in M$,
\[
\E{677}:\qquad
x=y\op\bigl(x\op((y\op x)\op y)\bigr).
\]
Prove that, for every $x\in M$,
\[
\E{255}:\qquad
x=((x\op x)\op x)\op x.
\]
\end{problem}

The following hints are meant to be used in order.  They are deliberately phrased as a
proof path, with the target of proving the identity $\E{255}$ from $\E{677}$ and
finiteness.

\begin{warning}
Do not assume associativity, commutativity, an identity element, inverses, right
cancellation, or quasigroup structure.  The only cancellation allowed in the hints below
is left cancellation after it has been proved.
\end{warning}

\section*{Hints}

\begin{enumerate}[label=\textbf{\arabic*.}]
\item \textbf{First obtain left division.}
Fix $a\in M$ and let $L_a(t)=a\op t$.  From $\E{677}$, with $y=a$, show that for every
$b\in M$,
\[
b=a\op\bigl(b\op((a\op b)\op a)\bigr).
\]
Thus $L_a$ is surjective, hence bijective because $M$ is finite.  Define
$a\backslash b$ to be the unique element $u$ with $a\op u=b$.  Then
\[
a\backslash b=b\op((a\op b)\op a).
\]
All later cancellations should be read as cancellations by one of these bijective left
translations.

\item \textbf{Derive the two basic identities.}
Show that, for all $y,z\in M$,
\[
z=(y\op z)\op\bigl((y\op(y\op z))\op y\bigr), \tag{T}
\]
and
\[
(y\op(y\op z))\op y
=z\op\bigl(((y\op z)\op z)\op(y\op z)\bigr). \tag{K}
\]
For (T), substitute $y\op z$ for $x$ in $\E{677}$ and cancel $y$.  For (K), compare
(T) with $\E{677}$ after replacing $y$ by $y\op z$, then cancel $y\op z$.

\item \textbf{Reduce the problem at a fixed point.}
Fix $x\in M$ and put
\[
a=x\op x,
\qquad
q=a\op x=((x\op x)\op x),
\qquad
p=x\backslash x.
\]
The desired identity at $x$ is simply
\[
q\op x=x. \tag{1}
\]
Prove the useful strengthening: if $u\op x=x$, then $u=q$.  Indeed, (T) gives
$x=x\op(x\op u)$, while $\E{677}$ at $(x,y)=(x,x)$ gives $x=x\op(x\op q)$; two left
cancellations give $u=q$.

Thus $\E{255}$ at $x$ is equivalent to the existence of a right fixed point over $x$.
It is also equivalent to each of
\[
q\backslash x=x,
\qquad
((x\op x)\op q)\op(x\op x)=x,
\qquad
(q\op x)\op q=p. \tag{2}
\]
Use (T) with $(y,z)=(x\op x,x)$, and the formula for $q\backslash x$, to verify these
equivalences.

\item \textbf{Introduce the fixed maps.}
For the same fixed $x$, define
\[
H(t)=(x\op t)\op x,
\qquad
F(t)=x\op(t\op x).
\]
Prove
\[
x\backslash t=t\op H(t),
\qquad
F(L_x(t))=L_x(H(t)). \tag{3}
\]
Hence $F$ and $H$ are conjugate self-maps of $M$.  Show that (1) is equivalent to
either of these maps having a fixed point.  If a fixed point exists, it is unique.

\item \textbf{Study the orbit of $x$ under $L_x$.}
Let
\[
c_0=x,
\qquad
c_{k+1}=x\op c_k.
\]
Since $L_x$ is a permutation of the finite set $M$, the sequence is periodic.  Let $d$ be
the least positive period of the orbit of $x$, and read all subscripts modulo $d$.  Put
\[
d_k=c_k\op x.
\]
Using $\E{677}$ with parameter $x$, prove the recurrence
\[
c_{k-1}=c_k\op d_{k+1}. \tag{4}
\]
Then show
\[
p=c_{d-1},
\qquad
q=d_1=c_{d-2},
\qquad
p\op c_1=q. \tag{5}
\]
A further useful fact is this: if $d_j=x$, then $j\equiv d-2\pmod d$.

\item \textbf{Eliminate the smallest periods.}
If $d=1$, then $x\op x=x$, so $q=x$ and (1) follows.  The case $d=2$ contradicts the
last statement in Hint 5.

For $d=3$, write
\[
c_0=x,
\qquad
c_1=a=q,
\qquad
c_2=p.
\]
Use (4) and (K) to prove successively
\[
x=a\op(p\op x),
\qquad
a\op p=x,
\qquad
p\op x=p,
\]
then
\[
(a\op a)\op a=p
\qquad\text{and}\qquad
(a\op a)\op a=x\op(a\op a).
\]
Cancel on the left by $x$ at the final step to get $a\op a=a$, contradicting the
displayed identity $(a\op a)\op a=p$.  Thus any unresolved point has orbit period
$d\ge4$.

\item \textbf{Use finiteness on the right.}
Assume, for the sake of contradiction, that $q\op x\ne x$.  By Hint 3, no element
$u\in M$ satisfies $u\op x=x$.  Hence $x$ is missing from the image of the map
$R_x(t)=t\op x$.  Since $M$ is finite, $R_x$ has a nontrivial fiber: there exist
$y\ne z$ and $e\in M$ such that
\[
y\op x=z\op x=e. \tag{6}
\]
Apply (T) twice, once with $(y,z)=(y,x)$ and once with $(y,z)=(z,x)$.  After left
cancelling $e$, obtain
\[
(y\op e)\op y=(z\op e)\op z=e\backslash x. \tag{7}
\]
Moreover, if $y\op e=z\op e$, then a further left cancellation gives $y=z$.  Therefore
the map $t\mapsto t\op e$ is injective on every such collision fiber.

\item \textbf{Convert a collision into the fixed point.}
The clean finishing lemma would be:
\[
\boxed{
y\op x=z\op x=e,\quad y\ne z
\quad\Longrightarrow\quad
e\op x=x.}
\tag{C}
\]
If (C) is proved, then the problem is solved.  Indeed, (6) gives $e\op x=x$, and Hint 3
then forces $e=q$, so $q\op x=x$.

Here is a route toward (C).  Define
\[
\Phi(u,v)=\bigl(u\op v,
(u\op(u\op v))\op u\bigr).
\]
If $\Phi(u,v)=(r,s)$, prove from (T) that
\[
r\op s=v. \tag{8}
\]
Thus a forward $\Phi$-orbit may be written as
\[
\Phi^n(u,v)=(\alpha_n,\alpha_{n+2}),
\qquad
\alpha_n\op\alpha_{n+2}=\alpha_{n+1}. \tag{9}
\]

For a collision (6), equation (7) gives a common element
\[
t=(y\op e)\op y=(z\op e)\op z,
\qquad
e\op t=x,
\]
and hence
\[
\Phi(y,x)=\Phi(z,x)=(e,t). \tag{10}
\]
It remains to prove $t=x$.  Since $M$ is finite, the forward $\Phi$-orbit of $(e,t)$
eventually repeats.  The intended last move is to pull the first repeated pair backward
through the recurrence (9) until it reaches the two distinct predecessors $(y,x)$ and
$(z,x)$.  If this pullback can be justified using only left cancellation, then $t=x$;
consequently $e\op x=x$, and the main identity follows.

\item \textbf{Optional orbit test for period four.}
As a smaller arena for the preceding ideas, suppose the $L_x$-orbit of $x$ has period
$4$.  Write
\[
c_0=x,
\qquad
c_1=a,
\qquad
c_2=q,
\qquad
c_3=p,
\qquad
r=q\op x,
\qquad
s=p\op x.
\]
Use the recurrence (4), together with (T) and (K), to prove
\[
x\op a=q,
\qquad
a\op x=q,
\qquad
a\op r=x,
\qquad
x\op q=p,
\qquad
x\op p=x,
\qquad
p\op a=q,
\]
and
\[
r\notin\{x,a,q\}. \tag{11}
\]
If $r=p$, show further that
\[
q\op a=q,
\qquad
s=(q\op q)\op q=a\op(q\op q),
\qquad
s\notin\{x,a,q,p\}. \tag{12}
\]
These identities make period four a good test case for the collision lemma (C): the
right tail is forced to leave the principal orbit unless (1) already holds.
\end{enumerate}

\begin{goal}
Prove the boxed collision lemma (C), or prove any one of the equivalent identities in
(1)--(2), using the finite repetitions supplied above.  Once $q\op x=x$ is obtained for
the fixed but arbitrary element $x$, the identity $\E{255}$ follows for all elements of
$M$.
\end{goal}

\end{document}
